% ---------------------------------------------------------------------------------------------------- 7
\section{Reproducing Part I step by step}\label{sec:repro1}

This section says what to do, in order, to obtain every number of Part~I from public data and models. Each step
gives the procedure and its parameters, the data it produces, and a check against the numbers reported above. The
file of the repository that implements each step is listed in Appendix~\ref{app:map}.

Two kinds of numbers appear. \emph{Measured} numbers (wall-clock speedups and the cost tables) need a GPU and will
differ slightly on other hardware. \emph{Replayed} numbers (selection, accuracy drops, fronts, re-splits,
certificates, load, bandit, robustness, theory) are computed on the CPU from the recorded traces and a cost table,
and are deterministic: given the same traces and cost tables they reproduce to the last printed digit. The traces
and cost tables we measured are committed, so every replayed number can be reproduced without a GPU by starting at
Step~R\ref{st:calib}.

\begin{figure}[ht]
  \centering
  \resizebox{\textwidth}{!}{%
  \begin{tikzpicture}[
      font=\small, >={Stealth[length=2.2mm]}, node distance=4mm and 6mm,
      gpu/.style={rectangle, rounded corners=2pt, draw=iitb, fill=iitb!12, align=center, text width=2.55cm,
                  minimum height=1.25cm, inner sep=3pt},
      cpu/.style={rectangle, rounded corners=2pt, draw=black!60, fill=black!3, align=center, text width=2.55cm,
                  minimum height=1.25cm, inner sep=3pt},
      dat/.style={rectangle, draw=okgreen, fill=okgreen!8, align=center, text width=2.35cm, minimum height=0.8cm,
                  font=\footnotesize, inner sep=2pt}]
    \node[gpu] (r1) {\textbf{R1} Fine-tune backbone + 6 exits};
    \node[gpu, right=of r1] (r2) {\textbf{R2} Record validation traces};
    \node[dat, below=of r2] (tr) {traces: conf, pred, probs per exit; tuning / held-out halves};
    \node[gpu, right=of r2] (r3) {\textbf{R3} Measure the per-layer cost $c(n)$, $h$};
    \node[dat, below=of r3] (ct) {cost table $c(n)$, head cost $h$};
    \node[cpu, below=13mm of tr] (r4) {\textbf{R4--R6} Calibration tables, policies, replay};
    \node[cpu, right=of r4] (r7) {\textbf{R7} Certified selection};
    \node[gpu, right=of r7] (r8) {\textbf{R8} Timed runs on the GPU};
    \node[cpu, below=of r4] (r9) {\textbf{R9--R12} Fronts, re-splits, certificates, load};
    \node[cpu, below=of r7] (r13) {\textbf{R13--R14} Bandit, robustness};
    \node[cpu, below=of r8] (r16) {\textbf{R15--R16} LLM tracks, theory checks};
    \draw[->] (r1) -- (r2);
    \draw[->] (r2) -- (tr);
    \draw[->] (r3) -- (ct);
    \draw[->] (tr) -- (r4);
    \draw[->] (ct) -- (r7);
    \draw[->] (r4) -- (r7);
    \draw[->] (r7) -- (r8);
    \draw[->] (r4) -- (r9);
    \draw[->] (r7) -- (r13);
    \draw[->] (r8) -- (r16);
    \node[anchor=north west, font=\footnotesize, align=left] at ($(r9.south west)+(0,-2mm)$)
      {\tikz\draw[draw=iitb,fill=iitb!12] (0,0) rectangle (0.3,0.2);~needs a GPU\quad
       \tikz\draw[draw=black!60,fill=black!3] (0,0) rectangle (0.3,0.2);~CPU only (deterministic replay)\quad
       \tikz\draw[draw=okgreen,fill=okgreen!8] (0,0) rectangle (0.3,0.2);~data committed to the repository};
  \end{tikzpicture}}
  \caption{The reproduction pipeline of Part~I. Every replayed result (grey) depends only on the committed traces and
  cost tables (green), so it can be reproduced on a CPU.}
  \label{fig:pipeline1}
\end{figure}

% ---------------------------------------------------------------------------------- R0
\subsection{Environment}
\begin{itemize}
  \item \textbf{Software.} Linux, Python 3.12.3, PyTorch 2.14.0 (CUDA 13.0), transformers 5.17.0, datasets 5.0.1,
  accelerate 1.15.0, numpy 2.5.3, scipy 1.18.1, pandas 3.0.6, scikit-learn 1.9.1 (isotonic regression), matplotlib
  3.11.2. Every package version we used is listed in Appendix~\ref{app:env}.
  \item \textbf{Hardware we used.} Server: 2\x{} Intel Xeon Silver 4216 (64 threads), 4\x{} RTX A6000 (48\,GB) and 2\x{} RTX 3080 Ti,
  NVIDIA driver 580.159.04; every server number in this report was measured on an RTX A6000. Laptop: GTX 1650 (4\,GB), 4 CPU threads. Any CUDA GPU can replace either; the measured
  speedups then change with it, the replayed numbers do not.
  \item \textbf{Data and models} (all public, from Hugging Face or torchvision): GLUE MNLI (\nolinkurl{validation_matched},
  9{,}815 examples) and QQP (\nolinkurl{validation}, 40{,}430) from \nolinkurl{nyu-mll/glue}; CIFAR-100 (test set, 10{,}000
  images) from torchvision; WikiText-103 (\nolinkurl{wikitext-103-raw-v1}) from \nolinkurl{Salesforce/wikitext}; backbones
  \nolinkurl{textattack/roberta-base-MNLI}, \nolinkurl{roberta-base}, \nolinkurl{google/vit-base-patch16-224-in21k},
  \nolinkurl{Qwen/Qwen2.5-1.5B}, \nolinkurl{meta-llama/Llama-3.1-8B-Instruct}, and the LayerSkip checkpoints
  \nolinkurl{facebook/layerskip-llama3.2-1B} and \nolinkurl{facebook/layerskip-llama3-8B}.
  \item \textbf{Conventions.} Speedup = full-model latency divided by the policy's latency at the same batch size.
  Drop = full-model accuracy minus policy accuracy, in points, on held-out data. Policies never see labels at decision
  time; labels are used only offline, to build calibration tables and to score.
\end{itemize}

% ---------------------------------------------------------------------------------- R1
\begin{step}{st:train}{Fine-tune each backbone jointly with its exits (GPU)}
\begin{enumerate}
  \item \textbf{Attach exits} after layers 2, 4, 6, 8, 10 and 12 of the 12-layer backbone; the exit after layer 12
  is the full model. A text exit is the RoBERTa classification head applied to the hidden state of the first token:
  dropout $\to$ Linear(768, 768) $\to$ tanh $\to$ dropout $\to$ Linear(768, classes). A ViT exit is LayerNorm
  $\to$ Linear(768, 100) on the class token.
  \item \textbf{Initialise.} MNLI: start from \texttt{textattack/roberta-base-MNLI} and copy its classifier into all
  six exit heads; its class order differs from GLUE's, so map checkpoint class $k$ to GLUE label $[2,1,0]_k$ (with
  the mapping the checkpoint scores 88.0\%, without it 28.5\%). QQP: start from \texttt{roberta-base} with randomly
  initialised heads. CIFAR-100: start from \texttt{google/vit-base-patch16-224-in21k}.
  \item \textbf{Loss.} The sum over exits of the cross-entropy of each exit, weighted by depth: the exit after layer
  $l$ has weight $l/\sum_{e} l_e$, i.e.\ $2/42, 4/42, \dots, 12/42$ (as in DeeBERT/PABEE).
  \item \textbf{Optimiser.} AdamW with two parameter groups (backbone and heads, each with its own learning rate),
  linear warm-up over the first warm-up fraction of steps followed by linear decay to zero, bf16 autocast, gradient
  norm clipped at 1.0. The training set is shuffled with seed 0. Text: premise and hypothesis (or the two questions)
  are tokenised as a pair, truncated to 128 tokens and padded to the longest sequence of the batch. Vision: images are
  resized from 32 to 224 pixels, augmented with pad-and-crop and horizontal flip, and normalised with mean = std = 0.5.
\end{enumerate}
\begin{center}\small
\begin{tabular}{@{}lccc@{}}
  \toprule
  Hyperparameter & MNLI & QQP & CIFAR-100 \\
  \midrule
  Initial weights & \makecell{\texttt{textattack/}\\\texttt{roberta-base-MNLI}} & \texttt{roberta-base} & \makecell{\texttt{google/vit-base-}\\\texttt{patch16-224-in21k}} \\
  Epochs / batch size & 2 / 64 & 2 / 64 & 4 / 128 \\
  Learning rate backbone / heads & $2\cdot10^{-5}$ / $2\cdot10^{-4}$ & $3\cdot10^{-5}$ / $2\cdot10^{-4}$ & $5\cdot10^{-5}$ / $10^{-3}$ \\
  Weight decay / warm-up fraction & 0.01 / 0.06 & 0.01 / 0.06 & 0.05 / 0.1 \\
  Dropout / label smoothing & 0.1 / 0 & 0.1 / 0 & 0 / 0.1 \\
  Max length / image size & 128 tokens & 128 tokens & 224 px \\
  Seed & 0 & 0 & 0 \\
  \bottomrule
\end{tabular}
\end{center}
\outp{a checkpoint per task (backbone + six heads) and a table of the accuracy of every exit.}
\chk{full-model (layer 12) validation accuracy 0.8777 (MNLI), 0.9106 (QQP), 0.9331 (CIFAR-100), Table~\ref{tab:models}.
One to three GPU-hours per model on an A6000.}
\end{step}

% ---------------------------------------------------------------------------------- R2
\begin{step}{st:traces}{Record the validation traces (GPU)}
\begin{enumerate}
  \item Run the whole validation set through the trained model once, in fp16 autocast, batch size 256, computing
  every exit (no early exit here).
  \item For every input $j$ and exit $e$ store the softmax vector $p_{j,e}$, the confidence
  $\text{conf}_{j,e}=\max_c p_{j,e,c}$ and the prediction $\text{pred}_{j,e}=\arg\max_c p_{j,e,c}$, plus the label.
  \item Split the validation set into two halves with a random permutation of seed 0: the first half of the
  permutation is the \emph{tuning} half (called ``evolve'' in the files), the rest is the \emph{held-out} half. All
  tuning, tables and selection use only the tuning half; every reported accuracy is on the held-out half.
\end{enumerate}
\outp{one trace table per task with columns \texttt{idx, label, split, conf\_e, pred\_e, p\_e\_c}.}
\chk{MNLI tuning half has about 4{,}900 inputs and QQP about 20{,}200. The traces we recorded are committed, so all
later CPU steps can start from them.}
\end{step}

% ---------------------------------------------------------------------------------- R3
\begin{step}{st:profile}{Measure the per-layer cost $c(n)$ and the head cost $h$ (GPU)}
\begin{enumerate}
  \item \textbf{Timing network.} Build the model of Step~R\ref{st:train} with random weights (\texttt{torch.manual\_seed(0)})
  in the dtype to be timed (fp32 on the GTX 1650). Timing does not depend on weight values. Make a pool of 512
  synthetic inputs of length $T\in\{16,32,128\}$: token ids uniform in $[5, 50000)$, first token 0, true length
  uniform in $[T/4, T]$ with the matching attention mask. (ViT: random $3\times224\times224$ images, 197 tokens.)
  \item \textbf{What one timed run does} (the same code path is used for the policies in Step~R\ref{st:time}): embed
  the batch; run the layers up to the next evaluated exit; apply that exit's head and softmax; copy the vector of
  maximum confidences to pinned host memory and synchronise (one host sync per evaluated exit); then continue.
  \item \textbf{Batch sizes} $B\in\{1,2,4,8,16,32,64,128\}$. For each $B$ time two plans: \emph{full} (only the last
  exit evaluated) and \emph{all heads} (all six exits evaluated, nobody leaves). Warm up once, then repeat each plan
  $r=\min(20,\max(3,\lfloor 20\,\text{s}/(0.02\,\text{s}\cdot B\,T/128)\rfloor))$ times and take the median; run
  the order full, heads, full, heads and average the two medians of each plan.
  \item \textbf{Derive} the per-layer cost $c(B) = t_{\text{full}}(B)/12$ and the head cost
  $h=\operatorname{median}_B\,(t_{\text{heads}}(B)-t_{\text{full}}(B))/5$. For reference fit the affine model
  $t_{\text{full}}(B)=12(a+bB)$ by least squares weighted by $1/t_{\text{full}}$ (relative error), and record
  $\rho=a/(bB)$ at $B=8$ and 64 and the largest relative error of the fit.
  \item \textbf{Cost table used by all later steps:} make $c(B)$ non-decreasing in $B$ with isotonic regression,
  interpolate linearly to every integer $n\in[0,256]$, extrapolate linearly beyond the last measured $B$, and set
  $c(0)=0$.
\end{enumerate}
\outp{one cost table per device, sequence length and model.}
\chk{Table~\ref{tab:cost} (GTX 1650): $c(1)\approx1.8$--2.3\,ms, $h\approx0.4$--0.5\,ms, affine error 41/31/10\% at
$T=16/32/128$; about 5 minutes per sequence length. The machine must be otherwise idle: the timed path synchronises
with the host after every exit, so CPU load from other jobs lowers the speedups and spreads them across rounds.}
\end{step}

% ---------------------------------------------------------------------------------- R4
\begin{step}{st:calib}{Build the calibration tables (CPU)}
For a set of tuning traces and an exit plan $S$ (the full plan has all six exits):
\begin{enumerate}
  \item For each evaluated exit $e$, cut its confidences into 20 bins at the quantiles of $\text{conf}_{\cdot,e}$
  (duplicate edges merged).
  \item In each bin $k$ and for every exit $e'$ of the plan estimate $P(\text{correct at } e' \mid \text{bin } k
  \text{ at } e)$ as $(\#\text{correct}+30\,\bar a_{e'})/(\#\text{inputs}+30)$, where $\bar a_{e'}$ is the overall
  accuracy of exit $e'$ (shrinkage with 30 pseudo-counts).
  \item Make each column non-decreasing in the bin index with isotonic regression weighted by the bin counts.
\end{enumerate}
The expected gain of input $j$ at exit $e$ is then $g_j(e\to e')=\max(0,\,P_{e'}(\text{bin}_j)-P_e(\text{bin}_j))$.
\end{step}

% ---------------------------------------------------------------------------------- R5
\begin{step}{st:policies}{Implement the policies and their candidate grids (CPU)}
Every policy is a function called at each evaluated exit (except the last) with the confidences of the $n$ active
inputs; it returns which inputs leave and whether the whole batch stops.
\begin{itemize}
  \item \textbf{Cost rule (ours).} With $F=12\,c(B)+h$, $\ell=L-l_e$ layers left and marginal cost
  $\Delta=c(n)-c(n-1)$ (for $n=1$ use $c(2)-c(1)$): input $j$ leaves if $g_j(e\to E)/B<\lambda\Delta\ell/F$, equation
  \eqref{eq:indiv}. If any inputs remain ($R$), the batch stops unless some later evaluated exit $e'$ has
  $\sum_{j\in R} g_j(e\to e')/B \ge \lambda[(l_{e'}-l_e)\,c(|R|)+(e'-e)\,h]/F$, equation~\eqref{eq:coll}.
  Grid: 40 values $\lambda\in\{10^{-3.5},\dots,10^{1.5}\}$, log-spaced.
  \item \textbf{Cost rule + exit plan.} Same rule on a plan $S$; the 31 plans are all non-empty subsets of the five
  intermediate exits, plus the last exit. Grid: every second $\lambda$ (20) \x{} 31 plans = 620 candidates. The
  calibration table is built on the plan's exits only.
  \item \textbf{Per-sample threshold:} input leaves if its confidence $>\tau$; 36 values $\tau\in[0.40,0.995]$ evenly
  spaced. \textbf{Threshold + exit plan:} every second $\tau$ (18) \x{} 31 plans = 558 candidates.
  \item \textbf{Majority vote:} the batch stops when more than half of the active inputs have confidence $>\tau$;
  same 36 values of $\tau$.
  \item \textbf{Unified exit:} every input leaves at one fixed intermediate exit (only that head is evaluated); 5
  candidates.
  \item \textbf{Random search:} a 3-parameter collective rule (per-sample threshold $\tau$; the batch also stops when
  the remaining share of the batch is below \texttt{frac\_cut} and their mean confidence exceeds \texttt{conf\_floor});
  150 candidates drawn with seed 0 from $\tau\sim U(0.5,1)$, $\texttt{frac\_cut}\sim U(0,0.5)$,
  $\texttt{conf\_floor}\sim U(0.3,0.95)$.
\end{itemize}
The counts above are for the six exits of the encoders. With $E$ exits the plans are the $2^{E-1}-1$ non-empty subsets
of the intermediate exits plus the last one; the LLM classifiers of Step~R\ref{st:llm} have $E=7$, so 63 plans, 1{,}260
candidates for the cost rule with a plan, 1{,}134 for the threshold with a plan and 6 unified exits.
\end{step}

% ---------------------------------------------------------------------------------- R6
\begin{step}{st:replay}{Replay a policy on recorded traces under a cost table (CPU)}
For each batch (a set of $B$ trace rows): walk through the evaluated exits of the plan in order; at each one add
(layers since the previous evaluated exit) \x{} $c(n)$ + $h$ to the batch's latency, where $n$ is the number of
active inputs, as in equation~\eqref{eq:latency}; at the last exit return the remaining predictions; otherwise call
the policy with the recorded confidences, return the recorded predictions of the inputs that leave and remove them.
The per-batch drop is the full model's accuracy on the batch minus the policy's; the speedup is $F$ divided by the
mean batch latency. This replay is exact for accuracy and is the ``cost-table'' speedup of Tables~\ref{tab:matched}
and~\ref{tab:resplit}.
\end{step}

% ---------------------------------------------------------------------------------- R7
\begin{step}{st:select}{Certified selection (CPU)}
For each task, batch size $B$ and method:
\begin{enumerate}
  \item Split the tuning half with a random permutation (seed 0) into a \emph{fit} part (40\%) and a \emph{cert}
  part (60\%). Cut each into disjoint random batches of size $B$ (permutation seeds $B$ and $B+1$).
  \item Build the calibration tables on the fit part. Replay every candidate on the fit batches.
  \item Order the candidates by mean drop (ascending), breaking ties by speedup (descending), and keep only those that
  are strictly faster than every candidate before them: this is the \emph{path} from safest to fastest.
  \item Walk along the path on the cert part. For each candidate compute the $K$ batch drops $Z_1,\dots,Z_K$, their
  mean $\bar Z$ and standard deviation $s$, and certify it if $(\alpha-\bar Z)/(s/\sqrt K) > t_{1-\delta,K-1}$ with
  $\alpha=0.01$ and $\delta=0.10$. Stop at the first candidate that is not certified.
  \item Return the fastest certified candidate (or none).
\end{enumerate}
The \emph{old protocol} of the first report, used for comparison in Table~\ref{tab:resplit}, pools fit and cert,
replays every candidate on it and takes the fastest one with empirical drop $\le\alpha$.
\outp{the certified parameters of every method at every $B$.}
\end{step}

% ---------------------------------------------------------------------------------- R8
\begin{step}{st:time}{Timed runs on the GPU}
\begin{enumerate}
  \item Take the held-out half and cut it into disjoint random batches of size $B$ (permutation seed $100+B$); time
  $n/B$ of them, with $n=1{,}920$ inputs at lengths 16 and 32 and $n=768$ at length 128 on the GTX 1650, and
  $n=2{,}048$ on the A6000. Each batch is mapped to $B$ random rows of the synthetic input pool (seed $B$).
  \item Policies timed: the full model, all heads with no exit (to measure the head overhead) and the certified
  policy of each method from Step~R\ref{st:select}, with calibration tables built on the fit part.
  \item Run one warm-up pass. Then, in each round, run every policy once over the same batches in a fixed order
  (interleaved), timing the whole pass with a host synchronisation before and after. The device executes exactly what
  Step~R\ref{st:profile} describes; after each host sync the confidences computed by the random-weight network are
  replaced by the trained model's recorded confidences for the same rows, the policy decides, and the leaving rows are
  removed from the hidden states with \texttt{index\_select}.
  \item Rounds: 5 at lengths 16 and 32 and 3 at length 128 (GTX 1650); 5 on the A6000. For each round compute the
  paired speedup (full time / policy time of the same round); report the median and the range over rounds.
  \item Confirm that the decisions taken on the device equal the replay of Step~R\ref{st:replay} for every batch,
  and compute the held-out drop by replaying the same policy on all held-out batches.
\end{enumerate}
\chk{Tables~\ref{tab:t04}--\ref{tab:t05} (GTX 1650, about 63 minutes in total) and
Tables~\ref{tab:t11}--\ref{tab:t12b} (A6000, about 28 minutes). Held-out drops must match exactly; speedups within a
few percent on the same GPU model, \emph{on an otherwise idle machine} (a loaded CPU spreads the rounds, because this
path synchronises after every exit).}
\end{step}

% ---------------------------------------------------------------------------------- R9
\begin{step}{st:front}{Accuracy--speed fronts and the hard quarter (CPU)}
\begin{enumerate}
  \item Run certified selection (Step~R\ref{st:select}) once per accuracy target
  $\alpha\in\{0.25, 0.5, 0.75, 1, 1.25, 1.5, 2, 3\}$ points, reusing one path per method and $B$.
  \item Score each selected policy once on the held-out half (disjoint batches, seed $100+B$) with the cost table:
  held-out drop and speedup. This gives Figure~\ref{fig:front}.
  \item Read each front at ``held-out drop $\le1$ point'': the fastest such point is the entry of
  Table~\ref{tab:matched}.
  \item Define the \emph{hard quarter} as the 25\% of held-out inputs with the lowest mean confidence over the five
  intermediate exits, and report the drop separately on it and on the rest at the one-point target
  (Table~\ref{tab:hard}, mean over $B=8,16,64$).
\end{enumerate}
\chk{Tables~\ref{tab:matched} and~\ref{tab:hard} exactly; about 5 minutes on 6 cores.}
\end{step}

% ---------------------------------------------------------------------------------- R10
\begin{step}{st:resplit}{Re-split validity check (CPU)}
Pool both halves of the traces. For each of 30 seeds $s$: permute with seed $1000+s$ and split into fit (20\%), cert
(30\%) and test (50\%); run certified selection (fit, cert) and the old protocol (fit + cert pooled); replay both
choices on disjoint test batches (seed $s+7$). Report the share of seeds whose test drop exceeds one point, the mean
drop and the mean speedup, at sequence length 16 with the GTX 1650 cost table, for $B=8$ and 64 and the three methods
of Table~\ref{tab:resplit}.
\chk{Table~\ref{tab:resplit} exactly (after the cache fix); about 8 minutes for both tasks on 6 cores. Build the
calibration table inside every re-split job from that job's fit part; never reuse a table across jobs
(Appendix~\ref{app:issues}).}
\end{step}

% ---------------------------------------------------------------------------------- R11
\begin{step}{st:certcheck}{Finite-sample and label-free certificates, and less calibration data (CPU)}
Repeat Step~R\ref{st:resplit} (30 seeds, 20/30/50\%) at $B=8,16,64,128$ with the A6000 cost table at length 128 (text)
and the A6000 ViT cost table (CIFAR-100), for five methods, with three certificates applied to the $K$ batch means of
the cert part: the $t$-test of Step~R\ref{st:select}; Hoeffding--Bentkus~\cite{rcps}; and the Waudby-Smith--Ramdas betting
bound~\cite{wsr}.
For the two finite-sample bounds map the batch drop from $[-1,1]$ to $[0,1]$. For the label-free version use the
disagreement with the full model's prediction as the risk (it upper-bounds the drop on every batch). For the data
experiment subsample both the fit and the cert parts to 10\%, 25\% and 50\%. Repeat MNLI with $\delta=5\%$.
\chk{Table~\ref{tab:cert}; about 1 h 47 min on 14 cores (3 h on 6 cores).}
\end{step}

% ---------------------------------------------------------------------------------- R12
\begin{step}{st:load}{Fluctuating load (CPU)}
\begin{enumerate}
  \item Generate an arrival trace: Poisson arrivals whose rate swings between light and heavy load (a sinusoid with a
  20\,s period plus bursts). A single server takes $\min(\text{queue},128)$ requests whenever it is free and serves a
  batch of size $B$ in the full model's time from the cost table. Record the sequence of batch sizes (3{,}000
  batches; 3 arrival seeds). All methods are replayed on the same sequence.
  \item For every method prepare three variants: certified once at $B=64$ (``once@64''), once at $B=8$
  (``once@8''), and re-certified per bucket $B\in\{1,2,4,\dots,128\}$ (each batch uses the parameters of the largest
  bucket $\le B$).
  \item Report the throughput speedup ($\sum$ full-model time / $\sum$ method time over all batches) and the drop over
  all served held-out inputs. Pass: tuned once, within budget, within 5\% of the best per-bucket baseline and ahead
  of every tuned-once baseline.
\end{enumerate}
\chk{Table~\ref{tab:load}; about 7 minutes per task.}
\end{step}

% ---------------------------------------------------------------------------------- R13
\begin{step}{st:bandit}{Label-free bandit (CPU)}
\begin{enumerate}
  \item Replace ``correct'' by ``agrees with the last exit'': keep, per exit $e$ and confidence bin $k$ (12 bins),
  the estimated rate $D[e,k,e']$ at which exit $e'$ disagrees with the last exit; the gain of going from $e$ to $e'$
  is $D[e,k,e]-D[e,k,e']$. No label is used.
  \item At batch $t$, with probability $\min(1, t^{-1/3})$ run a \emph{probe} batch to full depth (it still answers,
  at full cost) and update every table cell its inputs touch; otherwise run the cost rule with an optimistic
  (upper-confidence, $\beta=1$) gain, so that with little data it continues rather than stops. Use the $\lambda$ that
  the cost rule selected on the tuning half.
  \item Run $T=3{,}000$ batches of $B=64$ against an oracle that knows the true tables; record cumulative regret and
  fit its log--log slope. Use ten seeds (0--9) and three probe schedules: $\min(1,t^{-1/3})$, $\min(1,t^{-1/2})$, and a constant rate
  of 5\% of batches (probe probability $0.05\,t^{0}$); and a warm start
  from 20 labelled tuning batches. Baseline: a CeeBERT-style UCB1 over per-sample thresholds with reward = mean
  confidence $-$ $\mu\cdot$mean depth fraction.
\end{enumerate}
\chk{slopes 0.55--0.77 (Section~\ref{sec:bandit10}); about 22 minutes for ten seeds and six settings.}
\end{step}

% ---------------------------------------------------------------------------------- R14
\begin{step}{st:robust}{Slowdown inputs and the trimmed rule (CPU, plus one GPU step)}
\begin{enumerate}
  \item \textbf{Attack inputs.} ViT/CIFAR-100: a DeepSloth-style attack~\cite{sloth} (re-implemented): projected
  gradient descent in the $32\times32$ image space (before the resize to 224), $L_\infty$ budget $\varepsilon=8/255$,
  30 steps, loss $\sum_e \mathrm{CE}(\text{logits}_e,\text{uniform})$ so that no internal exit becomes confident, on
  500 images taken from the tuning half; record their traces as in Step~R\ref{st:traces} (GPU, once). Text tasks:
  use the hardest 2\% of held-out inputs as natural slowdown inputs.
  \item \textbf{Trimmed rule.} In the collective sum of~\eqref{eq:coll} drop the $m$ largest gains,
  $m\in\{0,B/16,B/8\}$.
  \item \textbf{Bound check.} For $k$ attackers, build 300 batches of $B=64$ with $B-k$ clean held-out inputs $C$ and
  $k$ attack inputs $K$; check $\operatorname{depth}(C\cup K;m)\le\operatorname{depth}(C;m-k)$ on every batch for
  every $k\le m$ (violations must be 0). Also record the victims' latency relative to the full model for
  $k\in\{0,1,2,4,8\}$ (16 and 32 in the final runs) and the same for majority vote, consensus and the threshold, all
  tuned by certified selection, on 4 cost profiles per task.
\end{enumerate}
\chk{0 violations in 25{,}200 batches; Section~\ref{sec:robust5}; about 7 minutes. The unit tests of the simulator and the bound
(19 tests) must pass.}
\end{step}

% ---------------------------------------------------------------------------------- R15
\begin{step}{st:llm}{LLM tracks (GPU for training and timing, CPU for the rest)}
\textbf{Batched classifiers} (Tables~\ref{tab:t15}--\ref{tab:t16}):
\begin{enumerate}
  \item Prompts. MNLI: \texttt{Premise: \{premise\}\textbackslash nHypothesis: \{hypothesis\}} followed by
  \texttt{\textbackslash nRelation of the hypothesis to the premise (entailment, neutral or contradiction):}. QQP:
  \texttt{Question 1: \{q1\}\textbackslash nQuestion 2: \{q2\}} followed by
  \texttt{\textbackslash nDo the two questions ask the same thing (yes or no)?}. Truncate to 128 tokens.
  \item Model. Freeze the base weights (bf16) and add LoRA~\cite{lora} (rank 16, $\alpha=32$) to the attention and MLP projections
  of every layer. Exits: Qwen2.5-1.5B after layers 4, 8, \dots, 28; Llama-3.1-8B after layers 8, 12, \dots, 32. Each
  exit is its own copy of the final norm followed by Linear($d$, classes), applied to the hidden state of the last
  prompt token.
  \item Train with the depth-weighted cross-entropy of Step~R\ref{st:train}: learning rate $2\cdot10^{-4}$ (LoRA)
  and $10^{-3}$ (heads); 8{,}000 steps at batch 16 for both models (128{,}000 training examples); seed 0. These are the values of the
  committed logs \file{logs/cls\_\{mnli,qqp\}\_\{qwen,llama\}.log}. The training script's defaults are 4{,}000
  steps at batch 32, and \file{scripts/reproduce\_gpu.sh} runs Qwen with those defaults, so pass
  \texttt{--steps 8000 --bs 16} for Qwen as well. Record traces as in
  Step~R\ref{st:traces}.
  \item Measure the cost table (Step~R\ref{st:profile}) and time the certified policies (Steps~R\ref{st:select},
  R\ref{st:time}) on the \emph{trained} weights with the real prompts padded to 128 tokens, in bf16. Check that the
  device's predictions agree with the traces (97--100\% at every exit).
\end{enumerate}
\textbf{Decoding} (Section~\ref{sec:llm}): train exit heads with LoRA (rank 16) or frozen (rank 0) on WikiText-103
next-token prediction (6{,}000 steps), record per-token traces, then, assuming a pure fixed cost ($b=0$), compare a
static exit (one exit for every batch, mixing adjacent exits at random to meet the budget exactly) with an adaptive
whole-batch stop (stop at the first exit where the batch's mean confidence reaches $\tau$), both tuned on the tuning
half to a one-point budget, for $B$ from 1 to 64.
\chk{Tables~\ref{tab:t15}--\ref{tab:t16}: Llama 60--63 minutes and Qwen 17--21 minutes of timing per task; adaptive
gain 12\% at $B=1$ falling to 0.2\% at $B=64$.}
\end{step}

% ---------------------------------------------------------------------------------- R16
\begin{step}{st:theory}{Theory checks (CPU)}
Enumerate small instances (a few inputs, 3--4 exits, monotone and non-monotone calibration tables, affine and floor
costs), compute the optimal policy by exact dynamic programming over (exit, active set), and compare it with
rules~\eqref{eq:indiv} and~\eqref{eq:coll}. Then replay the corrected per-input rule (test every later exit) against
the implemented one on the real traces (QQP, MNLI; both cost tables; $B=8$--128) and compare them at the certified
one-point target and at matched held-out drop.
\chk{210/210 and 17/190 (collective), 301 and 165 (per-input), 40 real-trace settings with $-1.4\%$ at matched drop
(Section~\ref{sec:theory}); about 10 seconds for the exact checks and 20 minutes for the real traces.}
\end{step}

\subsection{Checking a reproduction}
\begin{itemize}
  \item Every number printed in Part~I can be recomputed from the result files and compared with the text; we did
  this for every table cell (all match) and every number in the text (the ten that do not match the committed files are corrected and marked
  above).
  \item We regenerated every CPU-derived result file from the committed traces and cost tables on 25 November (72
  steps, 5\,h\,16\,min on 6 cores, CPU only): 92 of 102 files matched to the last printed digit. The other 10 differed
  only in cost-rule rows because of the cache bug of Appendix~\ref{app:issues}; after the fix, runs with 3 and 6
  workers give identical files.
  \item GPU numbers cannot be checked without a GPU; their raw logs are committed next to the result files.
\end{itemize}
